%
% Chapter 2 - Background and Problem Definition
%
\chapter{Background and Problem Definition} \label{cap:background}

This chapter describes the institutional context that motivates the EduCal project, details the concrete problems found in the current process used by ISEL's academic services, and states the objectives the project sets out to achieve.

\section{Problem Statement} \label{sec:problem-statement}

ISEL's academic services are responsible for planning and coordinating a large number of recurring institutional activities every academic year, among others application periods, enrolment windows, exam periods, and the various administrative deadlines associated with them. Each of these activities is not a single, isolated event: it typically recurs across the year in multiple occurrences, is owned by a specific department or responsible party, and must respect a set of institutional constraints, such as not falling on a weekend, a national holiday, or a school vacation period.

Today, this planning is done using spreadsheets (Excel). While spreadsheets are flexible and familiar, they are not designed to encode this kind of institutional logic, which creates several concrete problems:

\begin{itemize}
    \item \textbf{No structured relationship between events and their occurrences.} A spreadsheet row cannot natively express that one academic objective (e.g. ``Enrolment Period'') has several dated occurrences throughout the year, each with its own metadata; this has to be maintained by hand.
    \item \textbf{No automated year-to-year transition.} Each academic year, the whole calendar has to be manually rebuilt from the previous year's sheet, with every date recalculated and re-checked against weekends, holidays, and vacation periods.
    \item \textbf{No enforcement of business rules.} Spreadsheets do not validate minimum durations, minimum gaps between events, or conflicts with weekends and holidays; these rules exist only informally and are checked by hand, which makes them easy to miss.
    \item \textbf{Prone to filling errors and duplicated effort.} Replicating the network of objectives, responsible parties, and instructions by hand every year is exhaustive and error-prone: one incorrect date can propagate through the rest of the plan.
    \item \textbf{No institutional workflow metadata.} Classifying events by category, planning status, or responsible department is not natively supported, so this metadata, where it exists at all, is scattered across free-text cells or separate documents.
\end{itemize}

These limitations directly motivate EduCal, a tool that models academic planning as structured data instead of free-form spreadsheet rows, so that year-to-year transitions, rule checking, and institutional metadata can be handled by software rather than by manual, error-prone repetition.

\section{State of the Art} \label{sec:state-of-the-art}

Before designing a purpose-built system, it is worth asking whether an existing, general-purpose calendar tool could already solve this problem. The two most used products in this space are Google Calendar and Microsoft Outlook/Calendar, both of which are mature, widely adopted, and already familiar to most staff at an institution such as ISEL.

Both tools are built for scheduling individual meetings and personal appointments, sharing calendars between people, sending invites with RSVP tracking, and booking shared resources such as meeting rooms. Both also support recurring events, generally based on the iCalendar \texttt{RRULE} standard (e.g. ``every year on this date'', ``every weekday''), which lets a single event repeat automatically according to a fixed pattern.

However, neither tool addresses the specific institutional problem described in Section~\ref{sec:problem-statement}:

\begin{itemize}
    \item \textbf{No hierarchical event/occurrence model with institutional metadata.} A recurring event in these tools is just one event that repeats; there is no concept of an academic objective owning multiple independently-manageable occurrences with institution-specific metadata, which is the exact problem EduCal sets out to remove.
    \item \textbf{No academic-year migration.} Standard recurrence rules repeat an event on a fixed pattern; they cannot re-derive a whole calendar of interdependent events for a new academic year while checking each one against that year's specific holidays and vacation periods. There is no equivalent of EduCal's Migration Engine.
    \item \textbf{No configurable institutional business rules.} Neither tool can validate, at creation or migration, whether an event falls on a weekend, holiday, or vacation period, or keeps a minimum gap from another related event, the constraints central to EduCal's Business Rule Engine; users would have to check these manually.
\end{itemize}

In short, Google Calendar and Outlook are well suited to personal and team scheduling, but they operate at the level of individual events, not at the level of an institutional planning process with hierarchical objectives, yearly migration, and configurable rule checking. This gap is what justifies building EduCal as a dedicated system rather than adopting an existing calendar product.

\section{Objectives} \label{sec:objectives}

EduCal aims to replace this manual, spreadsheet-based effort with a web application. The main objectives are:

\begin{itemize}
    \item \textbf{Baseline Creation:} create an academic calendar from scratch, structured around events and their occurrences, as the foundation for the rest of the project's functionality.
    \item \textbf{Automation:} automatically migrate the academic calendar from the current year to the next, instead of rebuilding it by hand every year.
    \item \textbf{Application of Rules:} check compliance with configurable business rules, such as whether events fall on holidays, weekends, or outside school vacation periods, replacing the informal, manual checks used today.
    \item \textbf{Centralization:} unify the planning workflow in a single system that eliminates redundant spreadsheets and ensures data integrity and consistency across the institution.
\end{itemize}
